\chapter{校园生活}
\SourceParagraph{本节内容汇集了南赫学院学长学姐们对于校园生活的经验与感悟。大学生活中的许多问题，往往需要亲身经历后才能理解。希望大家能够认真阅读这些分享，从前人的探索与思考中获得启发，少走一些弯路，也更从容地适应新的学习与生活环境。}
\section{新生适应与校园经验}
\subsection{志愿、学习与生活适应}
\SourceLabel{一、前言}
\SourceParagraph{南大的大学生活最大的特点就是高度自由、极度散养。没有老师督促、没有固定学习生活节奏，很多同学刚入校都会陷入迷茫，不清楚该参与哪些社团、是否要参加学生工作、各类比赛和社交是否有意义、志愿活动和社会实践值不值得投入。在南大，想要快速适应生活、过得如鱼得水，核心能力从来不是埋头死读书，而是高效收集信息、精准取舍、拒绝无效内耗。本文结合大一全年真实体验，摒弃空洞套话，整理出适配南大新生的全方位实用生存干货。}
\SourceLabel{二、志愿活动：够用即可，功利与体验双向兼顾}
\SourceLabel{1.}
\SourceParagraph{在南大综测规则中，志愿时长达到 40 小时即为加分上限。除非真心热爱公益事业、计划长期深耕相关领域，否则无需盲目堆砌志愿时长。超额的志愿时长性价比极低，就像高中的三好学生荣誉一样，有则锦上添花，没有也不会影响评优、升学和综测核心成绩。}
\SourceLabel{2. 志愿选择技巧}
\SourceParagraph{尽量避开内容枯燥、纯苦力、无福利的零散志愿活动。优先选择校内社团、官方组织牵头开展的志愿项目，这类活动流程规范、体验感更好，既能轻松攒够所需时长，还能结识不同专业、三观契合的同学，拓宽自身的人脉和信息渠道。}
\SourceLabel{3. 优质志愿组织推荐}
\SourceParagraph{结合自身一年体验以及身边同学的真实反馈，真诚推荐两个高性价比公益组织，一是天健社，作为校内老牌优质社团，志愿活动形式丰富、氛围轻松，不搞形式主义，参与门槛低，非常适合大一新生入门；二是 AS 星空下的守望者组织，活动立意温暖、体验感极佳，福利完善，在积累志愿时长的同时，还能收获正向情绪价值和优质社交。大一会有百团大战，各个社团都会参加。大家也可以借此机会去了解更多的社团，去选择自己感兴趣的。}
\SourceLabel{三、社会实践}
\SourceParagraph{南大的社会实践没有强制内卷的要求，全校唯一的硬性要求，仅为毛概课程配套的实践任务，只需按要求完成课程任务即可。其余所有社会实践、调研走访、公益实践活动，完全根据个人兴趣自主选择。}
\SourceParagraph{多说一句，非常建议大一新生多参与各类社会实践和校园活动。大一课业压力小、空闲时间充裕，不用强求每一次参与都做出成果。大一参与各类活动的核心意义，在于拓宽眼界、积累人脉、收集各类信息、摸清自身兴趣和特长方向，通过不断试错找到自己愿意长期深耕的领域，为大二、大三的发展铺路。}
\SourceLabel{四、学习}
\SourceLabel{1. 放平心态，无需畏惧全英文教学}
\SourceParagraph{很多新生会对全英文授课课程感到焦虑，其实完全没有必要。南大各类课程的学习核心全靠自学，课堂讲授内容有限，考试重点、答题技巧、核心考点大多需要自主摸索，所有同学的起点基本一致，不用过度焦虑和盲目内卷。}
\SourceLabel{2. 技巧}
\SourceParagraph{这是亲身踩坑后总结的核心高分秘诀：一定要主动多和乐于助人、真诚靠谱的学霸交流沟通。很多课程的高分隐藏技巧，包括课程答案、期末出题范围、作业得分套路、考试避坑要点等。这些关键信息，大多是往届学长学姐传承、或者同学们自己找的。很多信息都是学姐丢分之后才知道的，后世定要引以为戒！}
\SourceLabel{五、生活与社交}
\SourceLabel{1. 宿舍关系：相敬如宾，不必强求成为挚友}
\SourceParagraph{一定要认清一个事实：大学舍友是随机分配的同住伙伴，并不是天生的朋友。无需强行和舍友深交、刻意磨合、勉强合群。三观不合、作息差异、相处不适都是大学宿舍常态。最优的宿舍相处模式是礼貌相处、互相尊重、互不打扰、安分守己。强行迁就他人、勉强融入，最终只会激化矛盾、消耗自身情绪和时间。真正的挚友是靠三观和爱好主动结交的，而非依靠宿舍随机分配。}
\SourceLabel{2. 社交避雷}
\SourceParagraph{大学生活中要重点警惕一类人：常年立“摆烂不学习”的人设，日常标榜自己全程摸鱼、从不复习，背地里却熬夜开灯偷偷刷题、默默内卷。这类人不仅会制造无效焦虑，还可能故意传递错误的学习信息、误导备考方向。建议在心里将其划为学习交流黑名单，无需深交、无需对峙、更不必当面质问。面对他们的言论淡然处之，听过即过，不较真、不内耗、不被其影响。}
\SourceLabel{3. 核心社交准则：极致减负，杜绝一切内耗}
\SourceParagraph{大学社交的终极奥义就是合得来就相处，合不来就疏远。不用刻意讨好任何人，不用勉强自己融入不属于自己的圈子。遇到消耗自己、相处不适的人，主动疏远、减少交集即可。如果遇到无法调和的宿舍矛盾，不用自我内耗、纠结忍耐，可以提前咨询辅导员、学长学姐换宿流程，大二可申请更换舍友，及时止损才是最优选择。}
\SourceLabel{4. 专注自我，坚持长期提升}
\SourceParagraph{摆脱无效社交和盲目跟风内卷后，一定要预留专属时间培养个人爱好、提升综合能力。大学最珍贵的就是充足的自由时间，深耕个人爱好、打磨专业技能、沉淀自我，远比跟风参加无用活动、勉强合群更有意义，也是大学生最核心的竞争力。}
\SourceBlankLine[2]
\subsection{信息渠道、学生工作与校园生活}
\SourceLabel{一、前言}
\SourceParagraph{珍惜鼓楼生活，积极探索，这里的探索包括但不限于鼓楼校区的地图，鼓楼周边的美食，南京的景点（玄武湖强推，晚上去放松身心，白天去疗愈自己），自己热爱的事情，信息来源的渠道，最好在大一结束之前确定自己未来的发展方向（本专业 or 跨专业，保研 or 出国，考公选调 or 升学等等），在大学里面最重要也最难的是不做 ddl 的奴隶，不做绩点的奴隶，不做欲望的奴隶，做你自己。忠告：他人得到的并不是你失去的。}
\SourceLabel{二、信息渠道}
\SourceParagraph{大学是一个巨大的 e 人养成器}
\SourceLabel{1. 人：}
\SourceParagraph{直系的学长学姐，大气学院的学长学姐，转群里其他院系的一些信息来源大手子（看他们发言的频率，建议私戳他们并保持礼貌的语言和诚恳的态度进行交流），书院辅导员，学院辅导员，专业课的老师，甚至社团的学长学姐，同年级的同学，敢于开口早晚能收集到信息的。}
\SourceLabel{2. 公众号：}
\SourceParagraph{南大公众号、有训书院、南大育教、南大青年、南大招生小蓝鲸（欢迎补充）}
\SourceLabel{3. 群聊：}
\SourceParagraph{26 转群 1084954435，南赫民间群 1040472048，南哪 26 年级大群 826811581，26 有训闲聊群 960123383 等等（欢迎补充）}
\SourceLabel{三、学生工作}
\SourceParagraph{(以下将在书院、学院、学校层面以及一些挂靠学校行政机构的社团的工作经验统称为学生工作经验）}
\SourceParagraph{不否认学生工作有利于刷脸提高存在感，进而在评优的时候提高竞争力，但纯卷学生工作经历不是很明智的选择}
\SourceLabel{0. 先修营（或者类似性质的先去学校的活动）：}
\SourceParagraph{如果想提前认识导员和有训书院管理机构的一些老师和学生辅导员，或者提前积累公众号供稿经验，建议报名。包括了一些参观活动，感想稿撰写，拍摄参观视频的机会，当时还提供了一些有意思的民俗活动。结营要小组合作完成一项团日活动策划案，还是比较锻炼人的能力的。}
\SourceLabel{1. 书院层面：}
\SourceParagraph{有训书院的学生工作体系比较特殊，称之为“团学班宿”，顾名思义，这个机构打通了书院团委、书院学生会、班级和宿舍，和别的书院“团学联”不完全一样，但性质类似。班级的班长团支书、心理委员、文体委员、学习委员、宣传委员会自动对应到团学班宿的对应部门，也就是成功进入“团学班宿”，在班级里就有对应职务。往年会在 9 月军训的时候招新，今年貌似定在开学典礼后。特点是和辅导员、书院教务主任打交道最多，算是大家能接触到的最早的，工作成果最容易看见的学生工作类型，例如新媒体中心负责公众号的各类美编，文体素拓部负责策划一些有训大一新生会参加的活动，青年组织部（班长团支书）负责党团材料填写，团组织生活以及团员大会的开展。在大一结束前夕会进行换届，然后竞选成功会进入为期一年的部长或主席任职期。因为本人曾经在青组部任职，相对来讲青组部体感会轻松一些。}
\SourceLabel{2. 学院层面：}
\SourceParagraph{南赫的团委和学生会大致会在 11 月份招新，大一学生除了写文编做美编这种线上的事情，以及偶尔来鼓楼开展的活动需要对接，其他基本没有什么事情要干。}
\SourceLabel{3. 学校层面：}
\SourceParagraph{大一下学期团学班宿会有挂职的通知，大一点同学可能经验不够多不一定能上，会有笔面试，一般是大二部长团主席团成员参与。}
\SourceLabel{4. 社团层面：}
\SourceParagraph{实际有精力参加的社团基本也就 2-3 个，比较推荐一个纯兴趣社团一个挂靠行政部门的社团（或者有传承有底蕴的大社团）。这里推荐招生志愿者协会，“南星梦想计划”的发起社团，挂靠南大招办，参与者依靠回校宣讲活动可以获得大量物料以及志愿时长，理事长平时除了转发推文基本无事，也就招生季相对较忙。}
\SourceLabel{三、学习}
\SourceParagraph{选择＞努力＞0}
\SourceLabel{1. 选个好课：}
\SourceParagraph{一门好课能让你的体验和成绩双丰收。如果你已经查看过生存手册的学分绩计算方法，那么应该明白南大是总评多一点分，学分绩就会高一点（不知道是不是好事），而不是有些学校 89 分折算 3.7，90 分折算成 4.0 这样梯度制。为此红黑榜（nju.ys.al）则成为了重要参考依据。这里建议大家如果学期结束总评出来，可以及时在 \url{https://table.nju.edu.cn/external-apps/7aded834-74a2-43cc-b515-fb8e01656ef2/login/?page_id=zI1D} 填写你的课程评价，造福之后的学弟学妹们。}
\SourceLabel{2. 一些策略：}
\SourceParagraph{开学进入退补选阶段以后，可以在课上在半夜稍微蹲一下，或许有惊喜。同时大一下学期的课程压力较大，如果你在大一上学期很不幸像笔者一样没有抽中自然通识人文通识科光美育，那么建议在大一上学期把《北欧文化与社会》这门课修掉，这样下学期起码不用再上《跨文化交际》能轻松点。}
\SourceLabel{3. 一些安慰：}
\SourceParagraph{如果没抢到大红甚至碰上了大黑榜，也不要气馁，这种需要你选的课程大概率不被纳入保研学分绩的计算里面，只对大额奖学金、出国有一些影响，补救方法建议查看另一位学长撰写的出国篇。}
\SourceLabel{4. 专业特色：}
\SourceParagraph{我们专业不需要考虑本科转专业、大类分流，天然会少一些烦恼，但是纯英文授课又弥补了这一部分。不过不需要担心，一些专有名词在课堂上面出现的频率很高，多看几遍多查几次就记住了。而相对的，词汇偏难那么语法会简单一些。而且英文教材尤其是微积分的教学顺序更符合我们思考的顺序，可读性（翻译成中文后）会强一点。资料欢迎进入南赫民间群自取。}
\SourceLabel{四、生活}
\SourceLabel{1. 调整你的心态}
\SourceParagraph{大学和高中最大的区别在于，高中大多数时候有一条相对明确的路径：上课、考试、排名、升学。进入大学以后，这条路突然分裂成了很多条，而且每一条看起来都很重要。有人从大一开始卷绩点，有人忙学生工作，有人提前进组，有人考语言成绩，也有人每天在鼓楼和新街口研究下一顿吃什么，思考哪里是春天看樱花最棒的地方，研究玄武湖哪个角度能拍到最美的落日。}
\SourceParagraph{但实际上大家选择的路径不同，你们之间并不构成彻底的竞争关系。我并不反对竞争，但是竞争不应当成为蚕食良知的借口，不应该成为破坏同学关系的导火索。}
\SourceParagraph{也不建议大家把朋友圈、空间、小红书、抖音当作大学生活的排行榜。你能看到的往往是别人筛选过后的高光时刻，看不到的是准备过程、运气成分以及试错成本，甚至是高光背后的痛苦。用别人的结果反复审判自己，是性价比最低的一种内耗。}
\SourceParagraph{同时笔者还希望教你怎么应对失败。大学里的很多失败都没有想象中严重。一次没选上学生组织，不影响下一次继续报名；一门课成绩不理想，也不代表你不适合这个专业；第一次和舍友、同学相处得不够顺利，也不代表你不擅长社交。大一允许迷茫，也允许试错。最重要的是把你不喜欢，不愿意长期坚持的事情找出来，筛选掉。}
\SourceParagraph{别人提前确定了未来方向，也不代表他的选择一定适合你。还是前面那句话，他人得到的，并不是你失去的。}
\SourceLabel{2. 控制你的欲望}
\SourceParagraph{刚进入大学时，大家通常会同时想要很多东西：想要高绩点，想进学生会，想参加比赛，想认识很多朋友，想谈恋爱，想出去玩等等等等，但是人的精力和金钱都是有限的。什么都想要，最后往往是样样通，样样松。}
\SourceParagraph{建议大家每个阶段只确定一到两个核心目标。比如大一上可以优先适应课程、熟悉校园、建立稳定社交圈；大一下再逐渐考虑竞赛、科研、语言考试。不是所有机会都必须立刻抓住，也不是所有活动错过一次就永远不会再来。}
\SourceParagraph{同时要警惕“因为别人参加了，所以我也应该参加”的跟风。别人报了比赛，你不一定要跟着报，别人周末到处聚餐旅行，你完全可以选择在宿舍睡觉或者打游戏。适合别人的生活模板，不一定能直接复制到你身上。}
\SourceParagraph{消费方面同理。鼓楼周边、新街口以及南京各种商圈的诱惑确实很多，外卖、奶茶、聚餐、演出、旅游和网购单看一次都不算贵，但叠在一起以后，月底嘛。。。}
\SourceParagraph{建议大一刚开始时简单记录一两个月的消费，不一定需要精确到每一分钱，但至少要知道自己的钱大概花在了哪里。吃饭、交通、日用品这些属于基本开销，社交娱乐和冲动消费则最好预留一个固定额度。不要为了维持所谓的合群，频繁参加自己并不想去的高消费聚餐，也不要因为不好意思拒绝，就替不熟的人垫钱、借钱或者参与金额较大的拼单。}
\SourceBlankLine[2]
\subsection{选课、活动与人际交往}
\SourceLabel{关于通识课：}
\SourceParagraph{个人建议选一个事少的老师（当然，事少的前提下，还是尽量选给分好一点的），尤其是阅读和美育。为什么这么说呢，因为给分是有可能红转黑的，但是事少的老师至少能让你空出时间来学其他的课。忙活一学期只有 85，肯定不如一学期啥也不干拿 85。}
\SourceParagraph{补充一下，大学的平时分可能存在一定的随机性，并不是说你平时做得多就一定会给你高分，也不是说你当了组长、干了很多活，就一定会给你更高的分数。还是要权衡一下，建议是把自己的时间花在更值得的事情上，比如学习专业课，或者其他你感兴趣的事。}
\SourceParagraph{体育：可以综合兴趣和给分进行选课，部分大热课程可能要在补选的时候抢（这点美育通识同理），虽然说一般情况下，只要你足够耐心，热门课程也是能抢到的，但还是建议先抢个给分好点的课保底。个人建议，零基础尽量不要选择球类运动（羽毛球是这样的，其他的我也不是很确定）。课上那点时间其实根本不够学的。}
\SourceLabel{关于志愿时长：}
\SourceParagraph{前面说的已经很完善了，我补充一点，如果你平时很喜欢参与各种大型集体活动，比如什么合唱比赛跳舞比赛之类的，那其实可以不用额外做志愿活动，大型活动给的志愿时长非常丰厚。当然，如果做志愿是你的个人爱好，那可以根据自己的情况来。}
\SourceLabel{关于校园活动：}
\SourceParagraph{虽然我本人非常热衷于参与各种活动，但是确实损失了一定的学习时间。如果你要从功利的角度看，参加活动的性价比确实比较低，但我认为，还是不能从完全功利的角度来看大学生活。当然，参加活动也是有很多好处的，比如说可以有机会做一些你感兴趣平时又不一定有机会做的事，可以结交到志同道合的朋友，可以锻炼自己的表达能力，等等，这些收获是很难被量化的。所以我的建议是：有选择地参加活动，选择你真正感兴趣的活动。}
\SourceParagraph{关于社团/校学生会/书院团学联/校团委/其他官方组织（仅根据我个人了解的内容来写）：}
\SourceListItem{\textcircled{1}呢喃大多数兴趣类社团活动时间应该相对自由，可以根据自己的兴趣来报。不报也行，很多社团有那种对外开放的大群，可以直接加群聊天；还有一些官方组织比如主持人队、乐团、合唱团，这种对水平一般有一定要求，需要面试，感兴趣的话可以关注他们的招新。其他三个主要是干活，组织活动什么的，每个部门职能不一样，可以自行了解。总的来说，各种组织最多加 2-3 个就可以了，太多了忙不过来。}
\SourceListItem{\textcircled{2}关于面试：紧张是正常的，其实面多了就麻了。面试其实就是一种自我推销，你要在较短的时间内展现出自己的价值。我的建议是选择自己最感兴趣的部门来面试，人在做自己感兴趣的事时，往往能表现的更自然，思维更流畅。尽量不要在短时间内面试太多部门，可以对一个部门进行更深度的了解，这样更有助于通过面试。不过，我说实话，这些面试可能根本代表不了什么，千万不要因为没有通过面试感到挫败，完全没必要。}
\SourceLabel{关于人际交往：}
\SourceParagraph{这个事情比较复杂，长话短说：\textcircled{1}注意观察分寸，学会反思是个好习惯。\textcircled{2}真诚和礼貌是必杀技。\textcircled{3}没必要太过功利。\textcircled{4}害人之心不可有，防人之心不可无。如果真的遇到了“有毒的关系”，及时跑路，必要时可以求助外界，比如找导员（导员不能帮你解决所有问题，但是有训的导员人都很好，至少可以帮你解决一部分问题，帮你减轻情绪压力）。}
\section{食物相关}

\SourceParagraph{笔者口味偏江浙一带，且较清淡，不能吃辣，在评价一些食物时可能含有较深的怨念。}
\SourceListItem{1 食堂}
\SourceLabel{1.1 鼓楼食堂}
\SourceParagraph{总体来说非常难吃，并随机出现饺子煮不熟、铁板饭把鸡蛋烤焦等彩蛋。这里仅指出非常离谱的一些。}
\SourceParagraph{鸭血粉丝汤的内脏在去年非常神奇的在刚开学的时候处理的还行，在之后就令人怀疑它是否被处理过，不知今年刚开学时是否能吃一阵。重要时间节点为大部分家长离开南大校园一周左右。}
\SourceParagraph{根茎类的蔬菜基本都处在刚好煮熟的状态。}
\SourceParagraph{除了明确标明为辣的菜，剩下的菜（点心类除外）有约 1/4 是辣的。}
\SourceParagraph{第三食堂的外包的铁锅菜（？还是炒菜）之类的不错，注意有些是辣的，有点油。}
\SourceParagraph{其实还是有几个能吃的东西的，请自行探索。}
\SourceLabel{1.2 仙林食堂}
\SourceParagraph{考虑到开学典礼在仙林，之后也会因各种活动需要在仙林吃饭，这里也介绍一下。}
\SourceParagraph{仙林食堂有一些广受好评的菜，但其中的佼佼者：第四食堂的大盘鸡拌面、第五食堂的黄焖羊肉面，它们是辣的。尽管如此，如果有不能吃辣的小朋友执意尝试，那么出门左拐步行 1min 有一家茶百道。}
\SourceParagraph{重点介绍一下九食堂，因为它离大气科学学院院楼非常近。开学典礼后的午餐是会发餐券的，但是这并不意味着校园卡不可以在仙林食堂用。餐券对应的都是盒饭（可能是三个菜都是辣的那种），且开学典礼结束里面都是人（这里指出一个事实：25 届开学典礼上没有任何形式的点名），所以其实直接去其他食堂可能是更好的选择。在开学典礼结束后有各种活动，如 25 届大气科学学院在开学典礼那天的下午有一个类似开放日的活动，而且仙林校区非常大（千万不要被鼓楼校区的占地面积给误导了），所以这里写一下仙林校园内通勤问题。大气科学学院的院楼离开学典礼的地方是有一定距离的，不建议步行，校内是有很多共享单车、助动车、滑板（这个可能不是很多？）的。另外，九食堂和大气科学学院院楼都在远东大道上，杜厦图书馆在南雍大道上，这三个地方和东大门在实际走的时候在感觉上是直行的（沿着主干道路走），但需要注意东大门并不是出门坐地铁的那个门，所以路痴要千万在来的路上记路，如果不幸发现自己在东大门附近（比如看到南大国际会议中心，南大博物馆之类的红色建筑），请右转前往南大门乘坐地铁。最后，如果真的不认路并且走到东大门时已近绝望，不要忘记你可以在手机上叫车（你可以在等车期间进入南大国际会议中心，里面很凉快；顺带一提进去有一个卖零食什么的地方，里面印着南大标识的手工酸奶相当不错），并不是怎么来就要怎么回去，所以也不要干在马路对面沿着地铁轨道寻找可以过马路的地方之类的事。最后的最后，以上故事非虚构。}
\SourceParagraph{第九食堂是外包的，不过鉴于刚重新装修过，里面的东西暂时应该不会变。轻食店的食材都挺新鲜的，但尽管是轻食，也要注意它里面看上去是辣的东西确实是辣的（不是非常辣），单独的彩椒除外。烧烤有一些烤的可能有些老，粉撒的很多，味道也很重。串串火锅非常不推荐，锅底在放肉之前上面就飘了一层白沫，并且在各个方面（食材品质、味道等）都毫无优点。江浙一带的同学应该会很喜欢“鸡鸣汤包”，里面的东西没有辣的，味道也都不错。帝王蟹什么的好像已经消失了。}
\SourceListItem{2 鼓楼附近}
\SourceLabel{2.1 适合堂食}
\SourceParagraph{这里的分类是随便分的，分在早餐的店大多全天都开。}
\SourceLabel{2.1.1 早餐}
\SourceParagraph{拿校外快递的那条路上早餐店都是 7:00 及以后开门的。味道都不错，但最前面那家（里面很大坐的地方的）偏咸、辣。徐州特色菜煎饼之规则怪谈：1）特色菜煎饼的一份其实是两人份（量非常非常大）。2）如果你的手机在它应该在的地方签到，食堂还没开，没有现金且感到饥饿，你或许会很高兴知道它是可以赊账的。}
\SourceLabel{2.1.2 午餐及晚餐}
\SourceParagraph{还是拿快递的那条路上两侧：新开的一家和快递站在同侧的（好像是面馆？）里面的面不错，可以让他不要放辣；石锅饭被诟病不正宗，但正因如此酱不是很辣，味道也不是很浓，比较对我口味，但牛肉很辣而且品质不行；吃鳗鱼饭什么的那家店，澳洲肥牛不是澳洲的，新西兰芝士不是新西兰的；吊龙面那家不错，但好像没有蔬菜。同一条路左转烧烤尚可，但老板对凉皮凉面不要放辣的理解和我不太一样，可能酱本来就是辣的。}
\begingroup
\small
\setlength{\tabcolsep}{4pt}
\renewcommand{\arraystretch}{1.18}
\begin{longtable}{>{\RaggedRight\arraybackslash}p{0.05\textwidth}>{\RaggedRight\arraybackslash}p{0.25\textwidth}>{\RaggedRight\arraybackslash}p{0.62\textwidth}}
\multicolumn{3}{l}{\textbf{堂食店铺信息（由原截图转录）}}\\[0.35em]
\toprule
编号 & 店铺与截图信息 & 原稿评价或截图文字 \\
\midrule
\endfirsthead
\toprule
编号 & 店铺与截图信息 & 原稿评价或截图文字 \\
\midrule
\endhead
\bottomrule
\endfoot
a & La Mia Casa 漫味意式小馆（上海路店）\newline 评分 4.8，约 122 条评价；南师大、意大利菜；上海路南秀村 14-1 号 & 这家如果双休日去一定要记得定位置（除非一个人去）。牛排和沙拉都是双人套餐里的那一种最好吃，蘑菇汤也很好喝。冰激淋很好吃。里面绝大部分东西都很好吃。但是，鲈鱼处理的不干净，而且腻，品质也一般（可能是因为鲈鱼在可食用鱼类中本来就一般）。 \\
\midrule
b & H617 汉堡\newline 评分 4.3，约 31 条评价；南大/南师大、西式快餐；青岛路 10 号 105 室 & 里面的东西从汉堡到热狗到各种薯条都很好吃，是快餐，老板人很好。牛肉品质相当不错。 \\
\midrule
c & 泥炉烧肉师（新街口北门桥店）\newline 评分 4.8，约 123 元/人；北门桥路 3-11 号（全季酒店对面） & 烤活鳗已经被切段放在架子上烤了尾巴还在跳。 \\
\midrule
d & A One Kitchen·南京精菜馆（晶丽酒店店）\newline 评分 4.8，约 184 元/人；南大/南师大、南京/江浙菜 & 如果不点套餐的话跟店员说是南大的学生，可以打折（好像是 88 折）。有河豚，里面的特色菜都好吃。打卡送的豆花和酸奶都非常好吃，且它们的好吃不在于免费，就是非常好吃。 \\
\midrule
e & 南京大牌档（鼓楼紫峰店）\newline 评分 4.2，约 71 元/人；紫峰购物广场 & 美龄粥不错，鱼偏咸，鸭血粉丝汤不错。 \\
\midrule
f & 成都娃娃花园餐厅 Garden Restaurant \& L...\newline 评分 4.7，约 101 元/人；南大/南师大、川菜馆；南秀村 18 号 & 是川菜，但可以让他做不辣的。酸菜鱼不错。 \\
\midrule
g & 汉口路新开的川菜馆（陈老师左边） & 汉口路新开的川菜馆非常美味，就在陈老师的左边。我和我舍友去吃了一下，非常夯。\newline 有好吃的不辣的菜吗？\newline 这个不知道，我们吃的都是略带辣味的。 \\
\midrule
h & 同一截图中的菜品与点单记录 & 一个小份的鱼，两个炒菜，我建议四个人吃。三个人没吃完。\newline 打包回来一个炒菜，三个人吃可以少点一个。\newline 看完菜单想了半天也不知道一个人去这个店怎么吃。然后看到有一个哥们点了一个辣子鸡加米饭。\newline 但是我感觉他会被辣死。\newline 椒麻鱼一定要让他少放麻椒。这个椒麻鱼的麻椒是我吃过最麻的一个。\newline 辣椒感觉还好。 \\
\end{longtable}
\endgroup
\SourceParagraph{一般。}
\SourceLabel{2.2 适合外卖}
\SourceLabel{2.2.1 早餐}
\begingroup
\small
\setlength{\tabcolsep}{4pt}
\renewcommand{\arraystretch}{1.18}
\begin{longtable}{>{\RaggedRight\arraybackslash}p{0.29\textwidth}>{\RaggedRight\arraybackslash}p{0.63\textwidth}}
\toprule
店铺 & 原截图信息 \\
\midrule
\endfirsthead
\toprule
店铺 & 原截图信息 \\
\midrule
\endhead
\bottomrule
\endfoot
福桔家庭厨房 & 评分 4.1，2084 条评价，约 86 元/人；私房菜，南大/南师大；截图显示 17:00 营业、距目的地约 190 米。 \\
\midrule
南京油条乌饭王（新街口店） & 评分 4.4，品牌已售 9000+；约 23 分钟、2.4 公里；起送约 10 元、免运费，截图标注“分量很大”；可选网红麻团蛋黄、油条麻团乌饭等。 \\
\midrule
小飞象面包公社（珠江路店） & 评分 4.7，已售 4 万+；京东秒送约 23 分钟；截图标注“当天现烤（只卖一天）”“总统黄油，自然醇香”“奶油碱水棒”。 \\
\end{longtable}
\endgroup
\SourceLabel{2.2.2 午餐或晚餐}
\SourceParagraph{a 强推}
\begingroup
\small
\setlength{\tabcolsep}{4pt}
\renewcommand{\arraystretch}{1.18}
\begin{longtable}{>{\RaggedRight\arraybackslash}p{0.29\textwidth}>{\RaggedRight\arraybackslash}p{0.63\textwidth}}
\toprule
店铺 & 原截图信息 \\
\midrule
蔡澜点心·粤菜（南京 IFCX 广场店） & 截图展示“经典拼盘”和“现包现烤”点心。 \\
\bottomrule
\end{longtable}
\endgroup
\SourceParagraph{个人最爱点的外卖，但不是南京特有的。}
\SourceParagraph{b 其他}
\begingroup
\small
\setlength{\tabcolsep}{4pt}
\renewcommand{\arraystretch}{1.18}
\begin{longtable}{>{\RaggedRight\arraybackslash}p{0.29\textwidth}>{\RaggedRight\arraybackslash}p{0.63\textwidth}}
\toprule
店铺 & 原截图信息 \\
\midrule
\endfirsthead
\toprule
店铺 & 原截图信息 \\
\midrule
\endhead
\bottomrule
\endfoot
小辣椒湖南米粉（珠江路店） & 评分 4.6，品牌已售 2 万+；约 27 分钟、1.9 公里；起送约 10 元、免运费，截图标注“量也很大”；可选杂酱干拌粉、干拌木耳肉丝、麻辣鸡粉等。 \\
\midrule
Ecovege 轻餐（新街口店） & 评分 4.7，品牌已售 90 万+；约 28 分钟、3 公里；起送约 8 元、运费约 0.2 元，截图标注“外酥里嫩”；可选深海金枪鱼蔬、菠菜虾仁能量、阳光番茄肥牛等。 \\
\midrule
柒宝酱 D9 街区店 & 截图展示金枪鱼土豆泥三明治 Tuna Potato Sandwich 1 人份，以及金枪鱼至味碗 Tuna Bowl 黑米杂粮饭。 \\
\midrule
無法现烤三明治（东大店） & 评分 4.8，品牌已售 1 万+；起送约 3 元、免运费，截图标注“分量也多”。 \\
\end{longtable}
\endgroup
\SourceParagraph{第一家是辣的，即使你在备注里强调不要辣，但味道不错。}
\section{关于金钱观和理财QwQ}
\SourceParagraph{本来这里是有一段很长的东西的。但是那是笔者大一写的，略显幼稚。决定后续不再列出，只做提醒。}
\SourceParagraph{不要借贷，不要赌博，投资不要上杠杆，投资需谨慎。}
\SourceParagraph{不要攀比，不要露富。}
\SourceParagraph{尽快通过实践或者专业的教学掌握正常的金钱观以及理财知识。}
\section{如何使用 ai 辅助学习}
\begin{CompactGuide}
\SourceParagraph{当代的大学生已经习惯将自身的学习生活与 AI 相结合，善用 AI 的能力已经成为个人的核心竞争力之一。这一章会给大家一些食用 AI 的建议，希望大家都能适应这个日新月异的时代。}
\subsection{认知}
\SourceParagraph{掌握使用 AI 的方式，要合理认识 AI。许多人对人工智能都会有焦虑，它的发展速度太快，在我们还不知道该做什么、该学什么的时候就要适应这个 AI 的时代，被迫调整我们的学习生活方式。不仅如此，AI 让我们的未来蒙上了一层不确定性的阴影，没有人可以预测 AI 会带领我们走向何方。需要认识到，大模型本身有着大量数学与工程的支撑。它的出现不也是一蹴而就的，而是经历了近百年来无数科学家前赴后继的探索。尽管人类对大模型还有很多没有理解的地方，但是我们不能停滞于仅仅将大模型视作“黑箱”的观点，现阶段它所展现出的能力与可能性应该被我们认可。不要害怕 AI 会取代你，而要这样思考：“我和 AI 合作，能完成什么以前做不到的事？”。}
\SourceParagraph{掌握使用 AI 的方式，要保持批判的思维。AI 有时效性，除非联网往往不能给出最新的信息。AI 还会“一本正经地胡说八道”（称为“幻觉”）。它给出的信息，尤其是事实、数据、文献引用，必须亲自核实来源。永远不要盲目相信 AI 的答案。更重要的是记住 AI 只是你的工具，使用工具的最终目的是提升自己。即使 AI 给出了令你满意的答案也不要停止，永远多想一想，怎么用 AI 得到更多更真实的收获。}
\subsection{使用}
\SourceParagraph{如果你觉得 AI 不好用，那大抵是你用的 AI 不够好。相较于有着各种限制的境内 AI，推荐有条件的人都去尝试使用国外 ai，例如 Gemini，chatgpt 等。对于那些在获取上述的资源有困难的人，可以尝试使用这样一个网页，askmanyai.cn。这是一个一站式 AI 大模型聚合平台，集合了 DeepSeek、ChatGPT、Claude、Gemini、Kimi、Grok、豆包等国内外多个顶尖 AI 大模型。它支持多个 AI 同时回答问题，你可以通过不同的 AI 模型比较和分析答案，从而得到最合适的信息。}
\SourceParagraph{如果你觉得 AI 不好用，那大抵是你用 AI 不够熟练。在 AI 时代，答案的获取成本极低，但提出一个好问题的能力变得前所未有的重要。学习如何精准、清晰、有深度地向 AI 提问（这通常被称为“提示工程”或 Prompt Engineering），这本身就是一项关键技能。好的问题能引出高质量的答案，而模糊的问题只会得到平庸的回复。更多的尝试用 AI 解决问题，你会慢慢掌握提问的要诀。}
\SourceParagraph{如果你觉得 AI 不好用，那大抵是你不知道用 AI 可以做到什么。有时候正是想象力限制了你对 AI 的使用。想到 AI 能干这件事，相信 AI 有可能干成这件事，AI 才可能做到、做好。一些可能的使用情景包括：用 ai 整理复习提纲、编程辅助、阅读论文、批量翻译、文档审核、对比变化、PPT 生成等。}
\SourceParagraph{如果你觉得 AI 不好用，那大抵是你不知道该用 AI 去探索些什么。下面有几点建议：}
\begin{HandbookBlock}
\SourceLabel{1. 了解 AI 的发展状况，评估它可能存在的边界}
\SourceLabel{2. 思考自己所在学科、工作对其的支持}
\SourceLabel{3. 尝试开发、使用基于 AI 技术的工具}
\SourceLabel{4. 思考、探索更多可能}
\SourceLabel{5. 变更自身的行为模式}
\end{HandbookBlock}
\subsection{学习/科研}
\SourceLabel{分享一些具体的我认为比较靠谱的经验：}
\SourceParagraph{在问答中学习。不论哪一种大模型，都有着极为丰富、全面的知识储备，且绝大多数情况下都会无条件的服从你的任务。遇到不会的东西就提问，对于一开始的回复不满意就让它进行修正，也可以让 AI 做提问者检验你是否真正理解了知识。我经常让 AI 生成五种不同的代码来完成同一个任务，AI 还是很好玩的。}
\SourceParagraph{阅读论文。可以试着使用 txyz.ai，帮你快速完成“预读”，翻译并提炼核心内容。}
\SourceParagraph{辅助写作。每种大模型擅长的写作场景都有所不同。都能用，用完记得自己读一遍、改一改。洋文可以试着使用 www.grammarly.com 润色。}
\SourceParagraph{辅助编程。用自然语言描述任务，让 AI 生成代码，然后要求它逐行解释代码。有报错就直接把完整的报错信息复制给 AI，它通常能快速定位问题并给出修改建议。}
\subsection{AI 工具使用教学和具体案例}

\SourceParagraph{随着人工智能的发展，AI 已经逐渐成为大学学习和科研中的重要工具。对于本科生而言，学习使用 AI 并不是简单地学会向聊天机器人提问，而是学会如何利用 AI 提高获取知识、解决问题和探索未知领域的效率。}

\SourceParagraph{目前常见的 ChatGPT、Claude、Gemini、DeepSeek 等工具，大多基于大语言模型（Large Language Model，LLM）。简单来说，大语言模型能够理解自然语言，并根据上下文生成文字、代码、图片分析结果等内容。相比传统搜索引擎，AI 不只是帮助我们寻找已有信息，还可以参与分析、总结和创造。}

\SourceParagraph{当然，我们也需要正确认识 AI 的能力。AI 并不是一个永远正确的“知识库”，它可能产生错误的信息，也可能生成看似合理但实际上不存在的内容。因此，AI 最适合作为学习和科研过程中的助手，而不是替代自己的思考。}


\subsubsection{一、常见 AI 工具选择}

\SourceParagraph{目前没有一个 AI 模型能够在所有方面都绝对领先。不同模型在推理能力、代码能力、长文本理解、多模态分析等方面各有特点。对于本科生而言，没有必要花费大量时间比较哪个模型“最强”，更重要的是根据自己的任务选择合适的工具。}

\SourceLabel{1. ChatGPT+Codex}

\SourceParagraph{如果只准备选择一个 AI 作为主力工具，我个人更推荐先尝试 ChatGPT。它的优势不一定是每一项能力都绝对领先，而是功能比较完整，文字、图片、文件、代码和数据分析等任务都能够处理，移动端和电脑端使用也比较方便。}

\SourceParagraph{对于我们而言，ChatGPT 可以辅助理解数学推导，帮助编写和修改 Python、MATLAB 代码，也可以用于分析实验结果、整理课程笔记和辅助阅读英文论文。}

\SourceParagraph{如果只是偶尔使用，免费版本已经能够满足很多需求；如果长期用于学习、科研或者项目开发，可以考虑购买会员。AI 服务更新较快，因此这里不固定推荐某一种充值方式。一般来说，优先选择官方网站或者官方应用中的订阅入口；如果官方支付存在困难，可以选择正规的礼品卡（如Apple的礼品卡）或者其他可靠渠道。对于学生而言，没有必要同时购买多个会员，一个自己熟悉并能够充分利用的 AI，通常比多个偶尔使用的工具更有价值。}

\SourceLabel{2. Claude}

\SourceParagraph{Claude 是 Anthropic 开发的大语言模型，它比较擅长处理长文本和复杂代码。相比日常问答，它更适合需要大量阅读和分析的任务，例如阅读长篇论文、分析大型代码项目、修改文章结构等。}

\SourceParagraph{在科研过程中，Claude 可以帮助我们快速完成论文预读。例如面对一篇陌生领域的论文，可以先让 AI 帮助梳理研究背景、解释方法流程、总结创新点，然后再结合自己的阅读深入理解。对于需要处理大量文字或者代码的任务，Claude 往往具有较好的表现。}

\SourceParagraph{Claude 的使用门槛相对较高，部分同学可能会遇到支付或者访问限制（比如ip需要非常纯净这点，若不纯净的话很容易被封号）。因此，如果目前主要需求是课程学习，没有必要专门购买；如果已经开始科研，并且经常需要阅读论文或者分析代码，可以考虑使用。}


\SourceLabel{3. Gemini}

\SourceParagraph{Gemini 是 Google 推出的人工智能模型，它在多模态任务方面具有一定优势，例如图片理解、图表分析以及结合 Google 生态完成资料整理。}

\SourceParagraph{但是现在Gemini总体不算好用，如果你想追求高端机制的模型，个人还是推荐ChatGPT/Claude}

\SourceLabel{4. DeepSeek、Kimi、通义千问等国内工具}

\SourceParagraph{DeepSeek、Kimi、通义千问、豆包和腾讯元宝等国内工具的优势是访问方便，中文使用体验较好，很多基础功能也不需要付费。日常学习问答、中文资料整理、简单代码和文字修改，用这些工具通常已经足够。}

\SourceParagraph{国内模型之间没有必要做过于绝对的排名。某个模型可能在数学题上表现很好，换到论文阅读或代码项目中却未必合适。比较实用的办法是保留一个主力模型，再准备一两个免费的备用模型，遇到重要问题时交叉询问。个人感觉Deepseek稍微好用一点，但直至今日，DeepSeek的多模态功能还是一个问题（即识别图像）}


\SourceLabel{5. 怎么选择}

\SourceParagraph{普通学习和综合任务可以先用 ChatGPT；需要阅读长论文、分析大型代码项目时可以尝试 Claude；经常使用 Google 文档或者需要处理图片资料时可以考虑 Gemini；希望免费、方便地完成中文问答和基础任务时，DeepSeek 等国内模型已经能够解决很多问题。}

\SourceParagraph{这个分类只是使用建议，并不是固定结论。真正适合自己的模型，需要通过实际任务不断尝试。很多时候，模型之间的差距并没有想象中那么大，能够清楚描述问题、判断答案是否可靠，才是使用 AI 最重要的能力。}


\subsubsection{二、AI 工具订阅和使用建议}

\SourceParagraph{AI 产品更新速度很快，模型名称、价格和套餐内容都会不断变化。因此，在使用 AI 服务时，不建议完全依赖网上流传的某个固定教程，而应该以官方网站显示的信息为准。对于本科生而言，最重要的问题不是“怎么买到最便宜的会员”，而是判断自己是否真的需要付费。}

\SourceParagraph{如果只是偶尔使用 AI 解决课程问题，免费版本通常已经能够满足需求；如果每天都会使用 AI 辅助学习、阅读论文、编写代码或者参与科研项目，那么购买一个稳定的会员通常是值得的。相比同时购买多个会员，一个自己熟悉并能够充分利用的工具往往更加有效。}

\SourceParagraph{订阅时优先选择官方渠道（比如有Apple设备的可以通过外区Apple id+购买礼品卡的方式订阅）。例如 ChatGPT、Claude、Gemini 。，可以选择正规的礼品卡或者可靠的支付渠道。如果在其他平台够买的话，需要仔细甄别，尽量不要购买来源不明的共享账号或者所谓低价会员，因为这不仅可能随时失效，也存在聊天记录、上传文件和个人信息泄露的风险。}

\SourceParagraph{同时需要注意，聊天会员和 API 通常是不同的服务。网页端会员主要用于日常对话，而 API 更多用于程序调用、自动化工具和 Agent。购买聊天会员并不代表拥有无限 API 使用额度。如果未来需要使用 API 开发工具，应当单独了解计费方式，并注意设置消费限制。}

\SourceParagraph{对于学生来说，AI 工具最重要的是稳定融入自己的学习流程，而不是不断追逐最新模型。一个能够长期使用、熟悉操作方式并且知道如何提问的 AI，通常比不断更换工具更加有价值。}
\subsubsection{三、Prompt、Context 与 AI 幻觉}

\SourceParagraph{很多同学第一次使用 AI 时，会发现得到的回答并没有想象中那么好。实际上，很多时候问题并不在于 AI 能力不足，而在于没有向 AI 清楚说明自己的需求。使用 AI 的过程，本质上也是学习如何描述问题的过程。}

\SourceParagraph{Prompt 通常翻译为“提示词”，简单来说就是我们提供给 AI 的指令。一个好的 Prompt 并不一定需要复杂的格式，而是需要让 AI 理解自己的背景、目标和要求。例如，同样是询问一个知识点：“解释一下 Navier-Stokes 方程。”如果增加背景信息：“我是一名大气科学本科生，已经学习过基础微积分和流体力学，希望从物理意义开始理解 Navier-Stokes 方程，请先解释它描述的问题，再介绍数学形式，最后给出一个简单例子。”通常能够得到更加符合需求的回答。}

\SourceParagraph{在实际使用中，可以尝试告诉 AI 自己是谁、正在学习什么、已经掌握哪些内容、希望解决什么问题。阅读论文时，可以说明自己的研究背景，并要求 AI 重点解释方法和创新点。提供的信息越准确，AI 越容易给出有帮助的回答。}

\SourceParagraph{Context（上下文）指的是 AI 当前能够理解的信息。使用 AI 时，一个非常重要的经验是：不要只告诉 AI 你遇到了什么问题，还要告诉它问题产生的背景。}

\SourceParagraph{例如询问代码问题时，只告诉 AI：“我的 Python 代码为什么报错？”通常很难得到有效帮助。如果同时提供完整代码、报错信息、运行环境、数据格式以及希望实现的功能，AI 才能够真正参与问题分析。对于大气数据处理任务，还应该说明数据来源、变量含义、时间范围和空间范围。例如 ERA5 数据中的温度、位势高度、降水等变量，其单位、维度和物理意义都有区别，AI 如果不了解这些背景，很容易给出形式正确但科学含义错误的代码。}

\SourceParagraph{Hallucination（AI 幻觉）是使用 AI 时必须注意的问题。AI 有时会生成看起来非常专业，但实际上错误的信息，例如不存在的论文、错误的软件接口、虚假的数据来源，甚至生成能够运行但分析方法错误的代码。}

\SourceParagraph{因此，对于论文引用、实验数据、专业结论以及科研分析结果，都需要自己进行验证。AI 可以帮助我们快速获得信息、整理思路，但不能代替自己的判断。尤其是在科研中，代码能够运行并不代表结果一定正确，图像能够绘制出来也不代表其中包含合理的物理意义。最终的判断仍然需要依靠自己的专业知识。}


\subsubsection{四、Agent、Skills 与 MCP}

\SourceParagraph{除了传统的聊天式 AI，近年来 Agent 也是人工智能发展的重要方向。普通 AI 的使用方式通常是“人提出问题，AI 返回答案”，而 Agent 更接近“人提出目标，AI 自动规划步骤，并调用工具完成任务”。}

\SourceParagraph{例如，在大气科研场景中，一个 AI Agent 可以根据研究目标自动读取 ERA5 数据，检查数据格式，调用 Python 完成数据处理，绘制温度、降水或者环流图，并根据结果生成初步分析报告。相比传统聊天模式，Agent 的重点不只是回答问题，而是能够参与完整任务流程。}

\SourceParagraph{当然，对于本科生而言，现在并不需要深入开发 Agent。更重要的是理解这种趋势：未来很多科研和工程工作，可能不会完全由人独立完成，而会变成人与 AI 协作完成复杂任务。能够提出合理目标、检查 AI 结果并调整任务流程，可能会成为一种重要能力。}


\SourceParagraph{与 Agent 相关的另一个概念是 Skills。简单来说，Skill 可以理解为 AI 的“技能包”。普通 Prompt 是每次使用时临时告诉 AI：“请帮我完成这个任务。”而 Skill 是提前将一套固定流程整理好，让 AI 可以重复调用。}

\SourceParagraph{例如，对于大气数据分析任务，可以设计一个数据处理 Skill，使 AI 按照固定流程读取 NetCDF 文件、检查变量信息、处理缺测值、完成数据清洗、绘制温度或者环流图，并输出分析结果。这样以后遇到类似任务时，就不需要每次重新描述完整流程。}

\SourceParagraph{对于本科生来说，目前不需要大量制作自己的 Skills，但可以关注已有的工具和工作流。未来随着 AI 发展，很多专业任务可能都会逐渐形成对应的 AI 技能。例如气象数据处理、论文阅读、代码分析、文献整理等，都可能拥有更加成熟的 AI 工作方式。}


\SourceParagraph{MCP（Model Context Protocol）则可以理解为让 AI 连接外部工具的一种协议。通过 MCP，AI 可以在获得授权后访问文件、数据库、代码环境或者其他软件工具，从而不再局限于聊天窗口中的文字交流。}

\SourceParagraph{例如，一个连接 Python 环境的 AI 可以直接运行代码并检查结果；一个连接文献管理工具的 AI 可以帮助整理论文；一个连接数据文件的 AI 可以辅助分析实验数据。}

\SourceParagraph{不过，工具能力越强，也意味着需要更加注意权限和安全问题。使用 Agent、MCP 或类似工具时，不应该随意开放所有文件和账号权限，也不要安装来源不明的插件。对于科研数据、个人信息以及未发表成果，更应该谨慎处理。}


\subsubsection{五、AI 辅助课程学习}

\SourceParagraph{AI 最适合在学习中扮演一个可以交流的助教，而不是简单的答案生成器。遇到不理解的知识点，可以让 AI 从不同角度解释；复习课程时，可以让 AI 帮助整理知识框架；完成习题后，可以让 AI 检查自己的推导过程。}

\SourceParagraph{例如，做数学题或者理论推导时，不推荐直接询问：“这道题怎么做？”因为这样得到的往往只是一个答案。更推荐告诉 AI：“这是我的解题思路，请指出哪里出现问题，并给出下一步提示，不要直接告诉我完整答案。”这种方式能够让 AI 参与学习过程，而不是简单替代自己的思考。}

\SourceParagraph{对于大气科学课程，也可以采用类似的方法。学习流体力学时，可以让 AI 解释方程中各项的物理意义；学习动力气象时，可以讨论地转平衡、位涡和尺度分析；学习天气学时，可以让 AI 根据天气图提出判读问题；学习大气探测时，可以让 AI 帮助整理仪器原理和误差来源。}

\SourceParagraph{不过，课程学习中仍然应该以老师的讲义、教材和课程要求为准。AI 对某个知识点的解释可能正确，但可能采用另一套符号体系或者不同的讲解方式。如果课程考试按照老师的体系进行，就需要让 AI 结合课程材料回答，而不是完全依赖通用知识。}

\SourceParagraph{一个比较推荐的学习流程是：先自己阅读教材和课程资料，记录不理解的问题；再使用 AI 进行解释和讨论；最后自己整理笔记并完成练习。这样既能够提高效率，也能够避免因为过度依赖 AI 而失去学习过程。}
\subsubsection{六、AI 辅助论文阅读和科研}

\SourceParagraph{对于本科生而言，AI 最大的价值之一，是降低科研入门的门槛。刚开始接触科研时，很多同学面对论文最大的困难并不是完全看不懂某一个公式，而是不知道这篇论文研究了什么问题、为什么采用这种方法，以及它在整个研究方向中处于什么位置。AI 可以帮助我们快速完成论文预读，建立对研究内容的整体认识。}

\SourceParagraph{阅读一篇陌生领域的论文时，可以先让 AI 帮助梳理研究背景、解释核心方法、总结实验流程和主要结论。例如可以询问：“这篇论文解决了什么问题？”、“作者为什么选择这种方法？”、“实验结果说明了什么？”、“这项工作的创新点在哪里？”这些问题能够帮助我们快速找到论文的重点。}

\SourceParagraph{对于大气科学方向而言，AI 在论文阅读中的作用尤其明显。很多论文涉及复杂的物理过程、数学模型和大量英文表达，AI 可以帮助解释专业术语，梳理研究思路，辅助理解图表。例如阅读一篇关于气候模式、遥感或者数值模拟的论文时，可以让 AI 帮助分析数据来源、模式设置、变量定义以及实验设计。}

\SourceParagraph{一个比较推荐的科研工作流是：论文搜索 → Zotero 管理 → AI 辅助阅读 → Obsidian 整理知识 → LaTeX 撰写论文。Zotero 可以帮助管理大量文献，避免论文和引用混乱；AI 可以帮助完成论文预读、方法解释和英文理解；Obsidian 可以帮助建立自己的知识体系；LaTeX 则适合完成论文、报告和学术文档排版。}

\SourceParagraph{不过，需要始终记住：“AI 阅读论文，不等于你阅读论文。”AI 可以帮助我们快速进入一个研究方向，但它可能忽略论文中的限制条件、实验细节以及作者隐藏的假设。真正判断一项研究是否可信，仍然需要自己阅读原文、检查数据和理解方法。}

\SourceParagraph{在论文写作过程中，AI 也可以发挥作用。例如帮助调整文章结构、润色英文表达、检查语句是否清晰、生成 LaTeX 表格或者协助修改图注。但是研究问题、实验设计、主要论证和最终结论必须由自己完成。特别是在科研中，不应该让 AI 根据模糊印象生成参考文献，也不应该上传未经允许的实验数据和未发表成果。}


\subsubsection{七、AI 辅助编程和数据分析}

\SourceParagraph{对于大气科学学生而言，编程能力会越来越重要。Python、MATLAB 等工具已经广泛应用于气象数据处理、数值模拟、数据可视化和科研分析。AI 的出现降低了编程学习的门槛，使得没有大量编程经验的同学也能够更快完成一些数据处理任务。}

\SourceParagraph{使用 AI 辅助编程时，一个比较推荐的方式是：先明确任务需求，再让 AI 提供代码，然后自己理解代码逻辑，运行测试并不断修改。AI 更像是一名编程助手，而不是完全替代程序员。}

\SourceLabel{例如：}

\begin{HandbookBlock}
\SourceParagraph{“帮我写一个 Python 程序读取 ERA5 数据，并绘制某区域温度变化。”}

\SourceParagraph{“解释每一步代码的作用，以及为什么这样处理数据。”}

\SourceParagraph{“如果变量改成降水或者位势高度，需要修改哪些地方？”}
\end{HandbookBlock}

\SourceParagraph{这样的使用方式不仅能够解决当前任务，也能够帮助自己学习编程思想。如果只是直接复制 AI 生成的代码，即使程序能够运行，也很难真正理解其中的数据处理逻辑。}

\SourceParagraph{大气数据处理中尤其需要注意这一点。例如 ERA5、CMIP、WRF 等数据通常具有复杂的数据结构，变量名称、单位、时间格式、空间坐标和数据维度都可能影响最终结果。AI 可以帮助编写 xarray、NumPy、Matplotlib 或 Cartopy 代码，但需要自己检查变量定义、单位转换和结果是否符合物理规律。}

\SourceParagraph{例如程序成功绘制了一张温度变化图，并不代表分析一定正确。需要进一步检查温度单位是否正确、时间平均方法是否合理、区域范围是否符合研究目的，以及结果是否符合基本物理规律。科研计算中，“代码运行成功”和“结果科学可靠”是两个完全不同的问题。}


\subsubsection{八、GitHub 与 AI 工具结合}

\SourceParagraph{GitHub 是科研和编程中非常重要的平台。很多开源代码、数据处理工具以及 AI 相关项目都会发布在 GitHub 上。对于本科生而言，学会使用 GitHub 不只是为了下载代码，更重要的是学会阅读项目结构、管理自己的代码以及参与开源协作。}

\SourceParagraph{使用 GitHub 时，首先应该学会阅读 README 文件。一个成熟的项目通常会说明项目用途、安装方法、依赖环境和使用示例。不要只看到项目名称符合需求就直接运行代码，而应该关注项目最近更新时间、Issue 讨论、许可证以及是否仍然有人维护。}

\SourceParagraph{AI 工具也经常与 GitHub 结合。例如使用 Coding Agent 时，可以让 AI 阅读一个 GitHub 项目，解释代码结构，帮助修改功能或者解决报错。但是在使用之前，应当理解项目内容和权限要求，不要随意运行来源不明的脚本，也不要把个人账号信息、实验数据或者未公开代码上传到公开仓库。}

\SourceParagraph{对于大气科学学生来说，可以逐渐积累自己的代码仓库。例如整理常用的数据读取程序、绘图模板、论文复现实验代码等。随着时间积累，这些代码会成为自己的科研工具库，而 AI 可以帮助不断完善和维护这些工具。}


\subsubsection{九、推荐的 AI 工作流}

\SourceParagraph{AI 真正提高效率的关键，不是使用某一个工具，而是形成适合自己的工作流程。不同任务应该选择不同工具，而不是所有事情都交给同一个聊天机器人完成。}

\SourceLabel{课程学习：}

\SourceParagraph{课程资料 → AI 解释 → 自己整理笔记 → AI 检查理解}

\SourceParagraph{学习新知识时，应该先阅读教材和老师讲义，明确自己不理解的问题，再让 AI 进行解释。如果完全依赖 AI 总结知识，很容易得到看似完整但缺少个人理解的内容。}

\SourceLabel{论文阅读：}

\SourceParagraph{论文搜索 → Zotero 保存 → AI 预读 → 自己阅读原文 → 整理笔记}

\SourceParagraph{AI 适合帮助降低阅读门槛，但真正重要的论文仍然需要自己阅读关键部分，例如研究问题、方法、主要图表和结论。}

\SourceLabel{数据分析：}

\SourceParagraph{明确科学问题 → 检查数据 → AI 辅助代码 → 运行验证 → 分析结果}

\SourceParagraph{数据分析中最容易出现的问题，是直接让 AI “帮我分析数据”。实际上，在得到结果之前，必须先明确数据来源、变量含义和研究目标。}

\SourceLabel{论文写作：}

\SourceParagraph{自己确定研究内容 → AI 优化表达 → 自己检查逻辑 → 完成最终版本}

\SourceParagraph{AI 可以帮助提高写作效率，但不能代替科研思考。文章真正有价值的部分，来自自己的研究问题、实验设计和分析过程。}


\subsubsection{十、AI 时代本科生应该培养什么}

\SourceParagraph{AI 降低了获取信息的成本，但提高了提出问题和判断信息的要求。未来真正具有竞争力的人，并不是完全依赖 AI 的人，而是能够利用 AI 放大自身能力的人。}

\SourceParagraph{首先，专业基础仍然非常重要。AI 可以帮助计算、整理和解释，但无法替代自己理解大气运动规律、数学模型和物理机制。对于大气科学学生而言，流体力学、热力学、动力气象、统计方法等基础知识仍然是判断 AI 结果是否可靠的重要依据。}

\SourceParagraph{其次，数据和编程能力会越来越重要。未来的大气科学研究会越来越依赖数据分析和计算方法。掌握 Python、MATLAB、数据处理和基本编程思想，可以让自己更有效地利用 AI，而不是只能被动接受 AI 给出的结果。}

\SourceParagraph{最后，需要培养与 AI 协作的能力。未来很多科研任务可能不再是一个人独立完成，而是人与 AI、人与计算工具共同完成。能够提出清晰的问题、拆解复杂任务、检查 AI 输出并不断优化流程，会成为重要能力。}


\subsubsection{十一、使用 AI 时需要注意的问题}

\SourceParagraph{AI 确实可以生成完整答案，但如果只是复制粘贴，你得到的只是一次任务完成，而不是能力提升。更推荐的方式是：先自己思考和尝试，再让 AI 提供解释、检查和补充，最后由自己总结方法。}

\SourceParagraph{在学习和科研中，合理使用 AI 是一种能力，但完全依赖 AI 代写作业、实验报告或者论文，不仅会降低自己的学习效果，也可能违反课程和科研规范。不同课程和科研团队对于 AI 的使用要求可能不同，使用前应该了解相关规定。}

\SourceParagraph{同时，也需要注意信息安全。不要随意向来源不明的平台上传个人信息、实验数据、未发表论文或者项目代码。对于科研内容，应当确认所使用工具的数据政策，避免因为使用 AI 而造成不必要的信息泄露。}

\SourceParagraph{最终，希望大家形成的工作方式不是让 AI 替代自己完成所有事情，而是：人提出目标，AI 提供辅助，人进行判断和修改，AI 进一步优化，最终由人完成决策。}

\SourceParagraph{AI 时代真正拉开差距的，可能不是谁最早购买了某个会员，而是谁能够把专业知识、编程能力和 AI 工具结合起来，完成以前难以独立完成的事情。与其担心 AI 会不会取代自己，不如尝试从一个具体问题开始，学习如何与 AI 一起完成更有价值的工作。}

\subsection{进阶篇}
\SourceParagraph{未来再写}
\SourceParagraph{和 AI 打交道与人类打交道有相似处也有不同点。相同的是都要注意保护隐私以及注意道德问题等。不同点则处于更含混的地带，此刻我难以言说评判。AI 向我们展示了一个崭新的世界，但是在 AI 之外仍有着更多东西存在。Ta 不是终点，而应当成为我们独立思考的开始。}
\SourceLabel{注：}
\SourceLabel{1，感谢 CAC}
\SourceListItem{2，这一篇章混用了 AI、大语言模型、人工智能的概念，具体它们间有什么区别和联系呢，你们未来会有机会学到的...}
\end{CompactGuide}
\newpage
\section{AI 时代：如何与变化共处}

\SourceParagraph{我们正处在一个特殊的时代。回顾人类的发展历程，工具始终在不断拓展人的能力：文字帮助我们保存和传递知识，计算器帮助我们处理复杂计算，互联网帮助我们连接世界。而如今，以大语言模型为代表的人工智能，正在进一步改变我们获取信息、学习知识和解决问题的方式。它不再只是一个简单的工具，而正在逐渐成为许多人学习、工作和创造过程中的重要伙伴。}

\SourceParagraph{面对 AI 的快速发展，焦虑是一种非常自然的情绪。很多人会思考：未来还有哪些工作属于人类？我们花费多年学习的知识是否仍然具有价值？如果 AI 可以完成越来越多的任务，人类自身的意义又在哪里？这些问题目前没有确定的答案，也没有人能够准确预测 AI 会把世界带向何方。但正因为未来充满不确定性，我们更需要培养面对变化的能力。}

\SourceParagraph{回顾历史，每一次重要技术变革都会带来类似的不安。工业革命改变了人类的生产方式，计算机改变了信息处理方式，互联网改变了人与世界连接的方式，而人工智能正在改变知识生产和知识应用的方式。技术确实会淘汰一部分旧的工作模式，但同时也会创造新的机会。真正重要的，并不是试图与 AI 竞争那些它更擅长完成的事情，而是思考：在人与 AI 的协作中，我们能够发挥怎样独特的价值。}

\SourceParagraph{未来优秀的人，或许并不是完全依靠自己完成所有事情的人，而是能够有效利用工具、提出高质量问题、快速学习新知识，并进行独立判断的人。AI 可以帮助我们搜索信息、整理资料、生成方案，甚至完成复杂的计算和分析，但它无法替代我们确定目标、理解价值、承担责任，也无法替代我们对于世界的真实体验和思考。}

\SourceParagraph{因此，学习使用 AI 并不意味着放弃自己的能力。恰恰相反，在 AI 能力不断增强的时代，我们更需要重视那些真正属于人的能力：扎实的专业知识、清晰的逻辑思维、良好的沟通能力、独立判断的能力，以及持续学习和适应变化的能力。AI 降低了获取信息的门槛，却提高了辨别信息、提出问题和进行创造的要求。}

\SourceParagraph{随着 AI 的发展，未来很多具体技能的价值可能会发生变化。过去，人们可能需要花费大量时间记忆信息、重复完成机械性的工作，而 AI 正在逐渐降低这些任务的成本。但这并不意味着学习知识失去了意义。恰恰相反，当信息获取变得越来越容易时，真正重要的是我们是否能够理解这些知识，判断信息的可靠性，并将已有知识应用于新的问题。}

\SourceParagraph{一个容易被忽视的事实是，AI 的能力越强，人类自身的基础能力反而越重要。因为 AI 可以快速生成答案，但它不知道什么问题值得被提出；它可以帮助我们完成任务，但它无法替代我们决定目标；它可以模拟各种观点，但最终仍然需要人类进行判断和选择。未来人与人之间的差距，或许不再主要体现在谁掌握更多固定的信息，而体现在谁能够提出更好的问题、学习得更快、判断得更准确，并利用工具创造更大的价值。}

\SourceParagraph{因此，大学阶段的学习并不只是为了掌握某一套未来可能过时的技能，而是在培养一种面对未知世界的能力。专业知识、数学基础、逻辑思维、沟通表达和科研能力，这些看似与 AI 没有直接关系的能力，实际上会成为我们理解和使用 AI 的基础。只有真正理解一个领域，我们才能知道如何向 AI 提出有价值的问题，也才能判断它给出的答案是否可靠。}

\SourceParagraph{同时，我们也应该警惕另一种误区：认为 AI 足够强大，所以人类不再需要努力学习。工具越强大，使用工具的人反而越需要具备更高的能力。不会使用计算器的人无法充分利用计算器，不理解计算原理的人也无法判断计算结果是否合理。同样，不具备独立思考能力的人，即使拥有最先进的 AI 工具，也很难真正发挥它的价值。}

\SourceParagraph{大学阶段或许正是我们学习如何与 AI 共处的重要时期。我们可以尝试使用 AI 辅助学习、参与科研、解决实际问题，也可以探索它带来的更多可能。但与此同时，请不要把 AI 当作逃避思考的捷径。一个优秀的 AI 使用者，并不是把所有事情都交给 AI 完成的人，而是知道什么时候应该借助 AI，什么时候应该依靠自己的人。}

\SourceParagraph{AI 向我们展示了一个崭新的世界，但在 AI 之外，仍然存在更多值得探索的东西。它不是终点，而应当成为我们重新认识自身能力、拓展思考边界的起点。希望大家能够拥抱变化，也保持独立；善用工具，也保持思考；在这个快速变化的时代里，找到属于自己的方向。}
